\documentclass{article}
\usepackage[utf8]{inputenc}
\usepackage[T1]{fontenc}
\usepackage{geometry}
\usepackage{graphicx}
\usepackage{float}
\usepackage{hyperref}
\usepackage[english,indonesian]{babel}

\title{Optimasi Pemrosesan Template LaTeX Menggunakan Pendekatan Berbasis AI}
\author{Academic Support Agent}
\date{\today}

\begin{document}
\maketitle

\begin{abstract}

\end{abstract}

\chapter{PENDAHULUAN}


\chapter{TINJAUAN PUSTAKA}


% Original content (commented out):
% % Vorlage für Praktikumsberichte
%
% main.tex
%
% Das ist das zentrale Dokument.


%persönliche Infos
\input{Berichtinfos}

%allgemeine Einstellungen
\input{Einstellungen}


\begin{document}

\include{Kapitel/Titelseite}
%include bindet andere .tex Dateien ein und hilft somit bei der Erstellung eines übersichtlichen Dokumentes

\pagenumbering{gobble}
\pagenumbering{Roman}
%Ab hier Seitenzahlen mit Römischen Ziffern

\tableofcontents
%erstellt das vollautomatische Inhaltsverzeichnis

\clearpage
\pagenumbering{arabic}
%ab hier Seitenzahlen in Arabischen Zahlen

\include{Kapitel/Einleitung}

\include{Kapitel/Praktikum}

\include{Kapitel/Fazit}

\listoffigures
%erstellt das Abbildungsverzeichnis
\pagebreak

\listoftables
%erstellt das Tabellenverzeichnis
\pagebreak

\printbibliography
%erstellt das Literaturverzeichnis

%Falls eine Erkärung benötigt wird
\include{Kapitel/Erklärung}


\end{document}
%beendet das Dokument



\end{document}